\section{Création du système}
\label{sec:systeme}

Dans le système final (figure~\ref{fig:systeme}), le programme reste dans la mémoire interne du
processeur et les données sont stockées dans le bloc \texttt{quadram} de la
partie~\ref{sec:quadram}. Le bloc de plus haut niveau, \texttt{systeme}, est écrit en VHDL : il
instancie le processeur Verilog en tant que \og component\fg{} ainsi que le bloc
\texttt{quadram}, et les relie par des signaux internes.

\begin{figure}[H]
\centering
\begin{tikzpicture}[font=\small, >=Stealth, line width=0.6pt,
    bloc/.style={draw, rounded corners=2pt, fill=gray!8},
    sous/.style={draw, align=center, fill=white, font=\footnotesize},
    ram/.style={draw, align=center, fill=blue!7, minimum width=2.4cm, minimum height=0.95cm,
                inner sep=2pt, font=\footnotesize},
    bus/.style={font=\footnotesize, above}]

  % processeur
  \node[bloc, minimum width=3.6cm, minimum height=6cm] (cpu) at (0,0) {};
  \node[anchor=north, align=center] at (cpu.north) {\textbf{CPU : riscvproc}\\Verilog};
  \node[sous, minimum width=2.9cm] at (0,0.4) {mémoire programme\\4096 $\times$ 32 bits\\
        (\texttt{\$readmemh})};
  \node[sous, minimum width=2.9cm] at (0,-1.7) {registres\\x0 à x31};

  % bloc quadram
  \node[bloc, minimum width=6.4cm, minimum height=6cm] (mem) at (8.6,0) {};
  \node[anchor=north, align=center] at (mem.north) {\textbf{MEM : quadram}\\VHDL};
  \node[sous, fill=orange!12, minimum width=1.5cm, minimum height=4.1cm] (dec) at (6.6,-0.55)
        {déco-\\deur\\2 vers 4};
  \foreach \i/\y in {0/1.15, 1/0, 2/-1.15, 3/-2.3} {
    \node[ram] (r\i) at (9.9,\y) {Sync\_Ram \i\\1024 $\times$ 32 bits};
    \draw[->] (dec.east |- r\i.west) -- node[above, font=\scriptsize] {cs(\i)} (r\i.west);
  }

  % bus entre le processeur et la mémoire
  \draw[->] (cpu.east |- 0,2.1)  -- node[bus] {\texttt{address[11:0]}}     (mem.west |- 0,2.1);
  \draw[->] (cpu.east |- 0,1.4)  -- node[bus] {\texttt{dataToWrite[31:0]}} (mem.west |- 0,1.4);
  \draw[->] (cpu.east |- 0,0.7)  -- node[bus] {\texttt{we}}                (mem.west |- 0,0.7);
  \draw[<-] (cpu.east |- 0,-0.6) -- node[bus] {\texttt{dataToRead[31:0]}}  (mem.west |- 0,-0.6);

  % horloge et reset
  \draw[<-] (cpu.west |- 0,1.0)  -- ++(-0.9,0) node[left] {\texttt{clock}};
  \draw[<-] (cpu.west |- 0,0.2)  -- ++(-0.9,0) node[left] {\texttt{reset}};
  \draw[<-] (mem.south) -- ++(0,-0.6) node[below] {\texttt{clock}};
\end{tikzpicture}
\caption{Système final : processeur et bloc \texttt{quadram}. Les bits 11--10 de l'adresse
attaquent le décodeur ; les bits 9--0, la donnée à écrire, \texttt{we} et l'horloge sont communs
aux quatre RAM, dont les sorties sont reliées (bus trois états).}
\label{fig:systeme}
\end{figure}

% ---------------------------------------------------------------------------
\subsection{Accès aux données dans la mémoire externe}
\label{sec:lwsw}

Les données n'étant plus dans \texttt{memory}, seules les instructions de manipulation des
données changent : \texttt{lw} et \texttt{sw} passent maintenant par le bus
(code~\ref{lst:lwsw}). \texttt{sw} place l'adresse, la donnée et \texttt{we} = 1 sur le bus ;
\texttt{lw} place l'adresse, puis recopie \texttt{dataToRead} dans le registre destination. Les
autres instructions ne sont pas modifiées.

\begin{lstlisting}[style=verilog, caption={Instructions \texttt{lw} et \texttt{sw} reliées au bloc \texttt{quadram}}, label={lst:lwsw}]
`LW: begin
  address   = regs[`RS1] + `I_IMM;   // address on the bus (falling edge)
  @(negedge clock);                  // rising edge in between: the quadram samples the address
  regs[`RD] = dataToRead;            // the word read is valid now
end
...
`SW: begin
  address     = regs[`RS1] + `S_IMM; // address and data on the bus
  dataToWrite = regs[`RS2];
  we          = 1;
  @(negedge clock);                  // rising edge in between: the quadram writes the word
  we          = 0;
end
\end{lstlisting}

% ---------------------------------------------------------------------------
\subsection{Synchronisation avec le bloc quadram}
\label{sec:synchro}

Le bloc \texttt{quadram} est synchrone : l'écriture a lieu sur le front montant de l'horloge, et
la lecture passe par le registre d'adresse \texttt{read\_address} de \texttt{Sync\_Ram} : la donnée
lue ne correspond à une nouvelle adresse qu'après le front montant qui a mémorisé cette adresse.
Le processeur travaillant sur le front descendant, un accès se déroule ainsi :
\begin{enumerate}
\item front descendant : le processeur place l'adresse sur le bus et, pour \texttt{sw}, la donnée
et \texttt{we} = 1 ;
\item front montant suivant : la RAM sélectionnée écrit la donnée, ou mémorise l'adresse de
lecture ;
\item front descendant suivant : \texttt{lw} recopie \texttt{dataToRead}, valide depuis le front
montant ; \texttt{sw} remet \texttt{we} à 0.
\end{enumerate}

L'attente est réalisée par un second \lstinline[style=vlinline]|@(negedge clock)| dans
l'instruction : \texttt{lw} et \texttt{sw} durent deux périodes, les autres instructions une
seule. Le tableau~\ref{tab:chrono} et la figure~\ref{fig:lw} montrent cette séquence pour
\texttt{test\_base} exécuté sur le système. À 140~ns, \texttt{lw x9, 4(x2)} place l'adresse 104,
mais \texttt{dataToRead} vaut encore $-7$, le mot 96 écrit juste avant ; la donnée 107
n'apparaît qu'après le front montant de 145~ns. Sans l'attente, \texttt{lw} lirait donc l'ancien
mot.

\begin{table}[H]
\centering
\begin{tabular}{rcll}
\toprule
Temps & Front & Instruction & Bus et registres \\
\midrule
100~ns & $\downarrow$ & \texttt{sw x3, 4(x2)}  & \texttt{address} = 104, \texttt{dataToWrite} = 107, \texttt{we} = 1 \\
105~ns & $\uparrow$   &                        & la RAM 0 écrit 107 au mot 104 \\
110~ns & $\downarrow$ &                        & \texttt{we} = 0 \\
120~ns & $\downarrow$ & \texttt{sw x1, -4(x2)} & \texttt{address} = 96, \texttt{dataToWrite} = $-7$, \texttt{we} = 1 \\
140~ns & $\downarrow$ & \texttt{lw x9, 4(x2)}  & \texttt{address} = 104, \texttt{dataToRead} = $-7$ (mot 96) \\
145~ns & $\uparrow$   &                        & adresse 104 mémorisée : \texttt{dataToRead} = 107 \\
150~ns & $\downarrow$ &                        & x9 = 107 \\
160~ns & $\downarrow$ & \texttt{lw x10, -4(x2)} & instruction suivante \\
\bottomrule
\end{tabular}
\caption{Accès à la mémoire pendant \texttt{test\_base} sur le système}
\label{tab:chrono}
\end{table}

\capture{p5-lw.png}{ModelSim, dossier \texttt{exo/p5}, programme \texttt{test\_base} : signaux
\texttt{clock}, \texttt{pc}, \texttt{instr}, \texttt{we}, \texttt{address}, \texttt{dataToWrite},
\texttt{dataToRead}, zoom de 95 à 165~ns}
{Écritures puis lecture dans le bloc \texttt{quadram} : la donnée lue n'est valide qu'après le
front montant}{fig:lw}

% ---------------------------------------------------------------------------
\subsection{Micro-architecture du processeur}

Les ports, l'horloge et le reset nécessaires à cet interfaçage ont été ajoutés à la
section~\ref{sec:microarchi}. Comme le processeur travaille sur le front descendant et la
mémoire sur le front montant, aucune autre modification n'a été nécessaire.

% ---------------------------------------------------------------------------
\subsection{Bloc système en VHDL}
\label{sec:blocsysteme}

Le fichier \texttt{systeme.vhd} (code~\ref{lst:systeme}) contient :
\begin{itemize}
\item la déclaration du \texttt{component} \texttt{riscvproc}, avec le même nom et les mêmes ports
que le module Verilog, écrits avec des types VHDL (\texttt{std\_logic} et
\texttt{std\_logic\_vector}) ;
\item les signaux internes qui forment le bus entre les deux blocs ;
\item l'instance du bloc \texttt{quadram} générique de la partie~\ref{sec:quadram}, avec
\texttt{Nbits\_sel} = 2, \texttt{Nbits\_addr} = 10 et \texttt{Nbits\_data} = 32 : quatre RAM de
1024 mots de 32 bits, soit 4096 mots adressés sur 12 bits, comme le bus d'adresse du processeur ;
\item les conversions de type : les ports de \texttt{quadram} sont de type \texttt{unsigned},
ceux du processeur de type \texttt{std\_logic\_vector}.
\end{itemize}

\begin{lstlisting}[style=vhdl, caption={Extrait de \texttt{systeme.vhd}}, label={lst:systeme}]
component riscvproc is
  port (
    clock       : in  std_logic;
    reset       : in  std_logic;
    address     : out std_logic_vector(11 downto 0);
    dataToWrite : out std_logic_vector(31 downto 0);
    dataToRead  : in  std_logic_vector(31 downto 0);
    we          : out std_logic
  );
end component riscvproc;
...
CPU: riscvproc
  port map (clock => clock, reset => reset, address => address,
            dataToWrite => dataToWrite, dataToRead => dataToRead, we => we);

-- 4 RAMs of 1024 words of 32 bits: address(11 downto 10) selects the RAM,
-- address(9 downto 0) is the word inside the RAM
MEM: entity work.quadram(RTL)
  generic map (Nbits_sel => 2, Nbits_addr => 10, Nbits_data => 32)
  port map (clock => clock, we => we,
            address => unsigned(address),
            datain  => unsigned(dataToWrite),
            dataout => dataOut);

-- the quadram works with unsigned, the processor with std_logic_vector
dataToRead <= std_logic_vector(dataOut);
\end{lstlisting}

Le testbench \texttt{systeme\_tb.vhd} génère l'horloge (période 10~ns) et le reset (à 1 jusqu'à
12~ns) ; un processus de surveillance arrête la simulation au bout de 100~$\mu$s si le programme
n'atteint jamais \texttt{ebreak}. Les blocs des parties 2 et 3 sont réutilisés sans copie : ils
sont compilés depuis leurs dossiers, dans la même bibliothèque \texttt{work} que le processeur
(code~\ref{lst:compil}).

\begin{lstlisting}[style=code, caption={Compilation et simulation du système dans ModelSim}, label={lst:compil}]
vlib work
vlog -sv riscvproc.sv
vcom -93 ../p2_q5/ramchip.vhd ../p3_q4/decoder.vhd ../p3_q4/quadram_for.vhd systeme.vhd systeme_tb.vhd
vsim -voptargs=+acc -onfinish stop work.systeme_tb
run -all
\end{lstlisting}

ModelSim associe le \texttt{component} au module Verilog de même nom présent dans la
bibliothèque. Si ce module n'a pas été compilé, ModelSim ne donne qu'un avertissement
(\og Component instance "CPU : riscvproc" is not bound\fg{}) et simule le système sans
processeur : c'est ce qui nous est arrivé au premier essai, et la simulation s'est alors arrêtée
sur le délai de surveillance de 100~$\mu$s.

% ---------------------------------------------------------------------------
\subsection{Validation}
\label{sec:validsysteme}

\paragraph{Programme instr2.}
À la fin du programme, x1 = 1 et x2 = 6, et le mot 65 se trouve maintenant dans la RAM 0 du bloc
\texttt{quadram} (\texttt{RAMS(0)/RAM/ram(65)} = 6, figure~\ref{fig:instr2}). La simulation se
termine à 80~ns au lieu de 70~ns, car l'instruction \texttt{sw} dure deux périodes.

\capture{p5-instr2.png}{ModelSim, dossier \texttt{exo/p5}, programme \texttt{instr2} : signaux
\texttt{clock}, \texttt{reset}, \texttt{pc}, \texttt{instr}, \texttt{we}, \texttt{address},
\texttt{dataToWrite}, \texttt{dataToRead}, de 0 à 80~ns}
{Système complet avec le programme \texttt{instr2}}{fig:instr2}

\paragraph{Accès aux quatre RAM.}
Le programme \texttt{test\_memoire} écrit 11, 22, 33 et 44 au même déplacement 5 dans chacune des
quatre RAM (mots 5, 1029, 2053 et 3077), puis relit ces mots ; il teste aussi un déplacement
négatif, la frontière entre les RAM 0 et 1 (mot 1023) et une lecture juste après une écriture à
la même adresse. Les résultats sont ceux attendus (tableau~\ref{tab:memoire}). Pendant les quatre
écritures, le signal \texttt{cs} du décodeur prend successivement les valeurs \texttt{0001},
\texttt{0010}, \texttt{0100} et \texttt{1000} (figure~\ref{fig:memoire}). La simulation se
termine à 350~ns : le programme compte 23 instructions, dont 11 \texttt{lw} ou \texttt{sw} de
deux périodes, soit 34 périodes à partir de la première instruction à 20~ns.

\begin{table}[H]
\centering
\begin{tabular}{lll}
\toprule
Instructions & Destination & Résultat \\
\midrule
4 \texttt{sw} au déplacement 5 de chaque RAM & mots 5, 1029, 2053, 3077 & 11, 22, 33, 44 \\
4 \texttt{lw} de ces mots & x20, x21, x22, x23 & 11, 22, 33, 44 \\
\texttt{add} des quatre valeurs & x24 & 110 \\
\texttt{lw x25, -5(x8)}, avec x8 = 3082 & x25 & 44 \\
\texttt{sw x24, -1(x5)}, avec x5 = 1024 & mot 1023 & 110 \\
\texttt{lw x26, 1023(x0)} & x26 & 110 \\
\bottomrule
\end{tabular}
\caption{Résultats de \texttt{test\_memoire} sur le système}
\label{tab:memoire}
\end{table}

\capture{p5-memoire.png}{ModelSim, dossier \texttt{exo/p5}, programme \texttt{test\_memoire} :
signaux \texttt{clock}, \texttt{pc}, \texttt{we}, \texttt{address}, \texttt{cs} (\texttt{/systeme\_tb/UUT/MEM/cs}),
\texttt{dataToWrite}, zoom de 80 à 180~ns}
{Les quatre écritures de \texttt{test\_memoire} : le décodeur sélectionne successivement les
quatre RAM}{fig:memoire}

Ce test a révélé une erreur dans notre \texttt{Sync\_Ram} de la partie~\ref{sec:ram} :
l'écriture n'était pas conditionnée par \texttt{cs} (\lstinline[style=vlinline]|if we = '1' then|).
Comme les quatre RAM partagent l'adresse, la donnée et \texttt{we}, chaque \texttt{sw} écrivait
dans les quatre RAM au même déplacement, et les quatre mots relus valaient tous 44 (x24 = 176 au
lieu de 110). La condition est devenue \lstinline[style=vlinline]|if we = '1' and cs = '1' then|.
Le testbench de \texttt{quadram} de la partie~\ref{sec:quadram} détecte aussi cette erreur
(messages d'erreur sur les RAM 0 à 2) lorsqu'il est exécuté avec la RAM non corrigée.

\paragraph{Programme test\_base.}
Les registres ont les mêmes valeurs qu'à la section~\ref{sec:microarchi}, et les mots 104 et 96
se trouvent maintenant dans la RAM 0 du bloc \texttt{quadram} (107 et $-7$). La simulation se
termine à 200~ns au lieu de 160~ns : les quatre accès à la mémoire durent chacun une période de
plus.

% ---------------------------------------------------------------------------
\subsection{Différences avec le premier programme}
\label{sec:instr}

\acompleter{comparer \texttt{instr2.txt} avec le premier programme de test (\texttt{instr.txt},
sur Madoc) : décoder les deux programmes instruction par instruction, en particulier les adresses
des \texttt{lw} et \texttt{sw}, et expliquer pourquoi le premier programme devait être modifié
pour fonctionner avec le système.}
